\documentclass[11pt]{article}
\usepackage{mathunal-base}
\mathunalfuente{serif}
\providecommand{\nota}[1]{\par\smallskip\noindent{\small\itshape\color{unalblue}#1}\par\smallskip}
\newlist{romanitos}{enumerate}{1}
\setlist[romanitos]{label=\roman*.,itemsep=1pt,topsep=2pt,leftmargin=1.8em,font=\normalfont}

\renewcommand{\examtitulo}{Supletorio Segundo Examen Parcial de Métodos Numéricos (25\%)}
\renewcommand{\examfecha}{29 de octubre de 2007}

\begin{document}

\examheader

\vspace{4pt}
\noindent\textit{Duración: 1 hora y 50 minutos.}

\vspace{6pt}
\noindent Nombre \dotfill\ Carné \dotfill\ Grupo \dotfill

\vspace{8pt}
\noindent\textbf{TODA RESPUESTA DEBE ESTAR DEBIDAMENTE JUSTIFICADA. NO SE ACEPTAN PREGUNTAS DURANTE EL DESARROLLO DE ESTE EXAMEN.}

\vspace{10pt}
\noindent\textbf{1. (25\%)} Determine valores de $A$, $B$ y $C$ que hagan que la fórmula de cuadratura:
\[
\int_{-1}^2 xf(x)\,dx\approx Af(-1)+Bf(1)+Cf(2)
\]
sea exacta para todos los polinomios de grado tan alto como sea posible. ¿Cuál es el grado de precisión de la fórmula obtenida?

\vspace{5mm}
\noindent\textbf{2. (25\%)} Una \textit{cercha cúbica} $S$ para una función $f(x)$ definida en $[-3,1]$ está dada por
\[
S(x)=\begin{cases}
-\dfrac23(x+3)^3-\dfrac12(x+3)^2+\dfrac{20}3(x+3)-7, & -3\leq x\leq -1\\[4pt]
a+b(x+1)+c(x+1)^2+d(x+1)^3, & -1\leq x\leq 1
\end{cases}
\]
Si $S''(1)=1$, encuentre $a$, $b$, $c$ y $d$. Escriba una tabla de interpolación para esta cercha.

\vspace{5mm}
\noindent\textbf{3. (20\%)} Si $f(x)=\ln(5-x)$ se aproxima en $[1,3]$ mediante el polinomio interpolante de grado menor igual que 3 empleando los nodos de Chebyshev, encuentre una cota \textit{razonable} para el valor absoluto del error en esta aproximación.

\nota{Ayuda: $|E_n(x)|=|f(x)-P_n(x)|\leq\dfrac{2(b-a)^{n+1}}{4^{n+1}(n+1)!}\displaystyle\max_{a\leq x\leq b}\left\{\left|f^{(n+1)}(x)\right|\right\}$.}

\vspace{5mm}
\noindent\textbf{4. (15\%)} Para los datos
\[
\begin{array}{c|c|c|c|c}
x_k & 2 & 5 & 8 & 12\\ \hline
f(x_k) & 1 & 2 & -3 & -8
\end{array}
\]
halle el polinomio interpolante de Lagrange y úselo para aproximar $f(6)$.

\vspace{5mm}
\noindent\textbf{5. (15\%)} Para los datos
\[
\begin{array}{c|c|c|c|c|c|c|c}
x_k & -3 & -2 & -1 & 0 & 1 & 3 & 5\\ \hline
y_k=f(x_k) & 4096 & 729 & 64 & 1 & 0 & 64 & 4096
\end{array}
\]
la tabla de diferencias divididas es
\[
\begin{array}{ccccccc}
-3 & 4096 & & & & & \\
-2 & 729 & -3367 & & & & \\
-1 & 64 & -665 & 1351 & & & \\
0 & 1 & -63 & 301 & -350 & & \\
1 & 0 & -1 & 31 & -90 & 65 & \\
3 & 64 & 32 & 11 & \alpha & 17 & -8\\
5 & 4096 & 2016 & 496 & 97 & 17 & 0 & 1
\end{array}
\]
Utilice esta tabla para encontrar:
\begin{romanitos}
\item (10\%) El polinomio interpolante $p(x)$ para $f(x)$ en los nodos $-2$, $-1$, $0$ y $1$.
\item (5\%) El valor de $\alpha$ en la tabla de diferencias divididas.
\end{romanitos}

\vfill
\rule{\linewidth}{0.4pt}\\[1pt]
{\scriptsize\textit{Transcripción digital --- MathUNAL}\hfill\textcolor{unalblue}{\textbf{1}}}

\end{document}
